\newpage

\section{Related Work}
\label{sec:02_related_work}

\subsection{Dense Retrieval}
Dense retrieval represents queries and documents as dense vectors and ranks documents through learned similarity functions. DPR~\cite{karpukhin2020dense} established a widely used dual-encoder formulation, learning query and passage representations from supervised question--passage pairs. Subsequent research improved both training and representation granularity: ANCE~\cite{xiong2021approximate} mined challenging negatives from a dynamically updated corpus index, while ColBERT~\cite{khattab2020colbert} retained token-level representations and computed fine-grained matching through late interaction. Retrieval-oriented pretraining also became an important direction. For example, coCondenser~\cite{gao2022unsupervised} introduced a corpus-level contrastive objective to prepare the passage embedding space before supervised training. Together, these approaches show that retrieval performance depends on the representation architecture, negative sampling, and pretraining objectives, providing essential controls when evaluating changes to representation geometry.

More recent research has expanded toward zero-shot retrieval and general-purpose embeddings built from large language models (LLMs). HyDE~\cite{gao2023precise} generated hypothetical documents from queries and encoded them to retrieve real documents without task-specific relevance labels. Wang et al.~\cite{wang2024improving} used LLM-generated synthetic data to train text embeddings through contrastive learning, evaluating generalization on BEIR and MTEB. The Qwen3 Embedding technical report~\cite{zhang2025qwen3} further combined foundation models, synthetic data, and multi-stage training for multilingual text embedding and retrieval. These developments broaden the available retrieval backbones and training resources, while motivating evaluation beyond a single domain. For structured content, SANTA~\cite{li2023structure} introduced Structured Data Alignment and Masked Entity Prediction, demonstrating their effectiveness in code and product search. Building on this direction, the present study examines whether hyperbolic alignment within SDA improves S1000D XML retrieval and preserves relevant hierarchical relationships, with encoder and training conditions controlled to distinguish geometric effects from other sources of improvement.

\subsection{Hyperbolic Representation Learning}
The use of hyperbolic geometry to represent relational structure predates modern neural embedding methods. Krioukov et al.~\cite{krioukov2010hyperbolic} connected negative curvature with degree heterogeneity and clustering in complex networks, providing a basis for relating geometry to structural organization. In 2017, Nickel and Kiela~\cite{nickel2017poincare} introduced Poincar\'e embeddings for learning hierarchy and similarity jointly. On the WordNet noun hierarchy, they compared Poincar\'e, Euclidean, and translational embeddings across dimensions using reconstruction and held-out link prediction, with mean rank and mean average precision showing strong low-dimensional performance. Their 2018 study~\cite{nickel2018learning} investigated Lorentz-based optimization on WordNet, EuroVoc, ACM, and MeSH, evaluating relation ranking and correlations between embedding norms and taxonomic depth. Its improvements concerned the learned embeddings and optimization procedure; the Poincar\'e and Lorentz models describe the same hyperbolic geometry at matching dimension and curvature. Also in 2018, Ganea et al.~\cite{ganea2018hyperbolic} extended hyperbolic learning to classifiers and neural layers. Comparisons on binary SNLI entailment and synthetic noisy-prefix recognition showed clearer advantages for more tree-like prefix relationships, while SNLI results were broadly comparable to Euclidean counterparts.

Subsequent applications examined how these geometric properties interact with task objectives and model components. Hyper-Know~\cite{ma2021knowledge} combined Poincar\'e representations, knowledge-graph attention, and adaptive regularization for recommendation. On Amazon-book, Last-FM, and Yelp2018, held-out interaction ranking and component ablations assessed Recall@K and NDCG@K, reporting improvements over the evaluated baselines. These results support the combined approach, rather than attributing all gains to geometry alone. Weber et al.~\cite{weber2024flattening} trained segmentation models on Cityscapes and tested them on five additional driving datasets, comparing hierarchical, flat Euclidean, and hyperbolic classifiers. Flat Euclidean models remained strong accuracy baselines, whereas hyperbolic representations showed particular benefits for parent-level calibration under domain shift. More directly related to retrieval, HypRAG~\cite{madhu2026hyprag} evaluated Lorentz-based text encoders on MTEB and RAGBench, comparing embedding quality and generated-answer grounding and relevance. Pooling ablations and analyses of radial separation between general and specific concepts examined how document aggregation retained hierarchical information. Together, these studies motivate evaluating the modified SANTA model on S1000D XML retrieval through controlled comparisons of geometry and alignment components, alongside direct checks of the structural relationships retained by its representations. Retrieval gains alone would not establish successful structural learning.

